\chapter{Gestión de la Misión y del Proyecto}

\section{Manejo Sistémico e Hitos (Milestones)}
El diseño y lanzamiento de una plataforma en formato CubeSat 2U como INTISAT se articula formalmente bajo los ciclos de desarrollo internacional V-Model de ingeniería de sistemas espaciales impulsados por INRAS/PUCP. La trayectoria traza un puente estricto desde SRR (Requimientos de Sistema), superando la presente CDR (Critical Design Review) del modelo de ingeniería y finalizando la madurez de hardware comercial COTS en la AIT de su inminente Vuelo Final (Flight Model).

\section{Cronograma, WBS y Manejo Acoplado de Riesgos}
\begin{itemize}
    \item \textbf{Árbol Estructural Desglosado de Trabajo (WBS):} Subdivide y paraleliza los requerimientos operacionales y cargas de manufactura en un esquema matricial jerárquico. Las ingenierías ramificadas como el *Software del OBC*, *Manufactura Térmica EPS* y *Entrenamiento IA del Payload* progresan encapsuladas y confluyen transversalmente en hitos mecánicos vinculantes (eg., Integración de Mazo de Cables y pruebas \textit{FlatSat}).
    
    \item \textbf{Matriz de Control y Mitigación de Riesgos Técnicos (Risk Management):} Previsión estocástica atada al control de fallas del hardware. Tipifica desastres potenciales en órbita (como cortocircuitos causados por radiación o falla en soldaduras BGA tras vibración de cohete). 
\end{itemize}

Delinea además respuestas técnicas pre-vuelo ante infortunios inevitables: por ejemplo, ante sub-voltaje severo por eclipse persistente, la planificación delega el diseño inteligente preventivo (\textit{By Design}) la reacción de contingencia hacia el \textit{Modo de Seguridad (Safe Mode)}, deteniendo la degradación fatal, desconectando instrumentos colaterales en conjunto y enviando peticiones de emergencia a la red terrestre global a espera del \textit{Troubleshooting} en consola.
